% ============================================================
% Sección 3
% ============================================================
\section{3 EDO Autónoma. $y'(x)=g(y)$}

\begin{tcolorbox}[enunciado]
Las EDO autónomas son un caso particular de las separables y se pueden resolver de manera análoga.

\vspace{0.2cm}
\textbf{Clase:}
\begin{itemize}[leftmargin=*]
    \item Análisis Cualitativo asíntotas y Sol Analítica, EDO Autónoma 7 Sep 2026 MCA4
\end{itemize}

\textbf{Código:}
\begin{itemize}[leftmargin=*]
    \item Github: \texttt{1.3.Sol\_Analitica\_Bifurcacion\_Node-Saddle\_dx\_mu-x2\_mu\_neg.py}
\end{itemize}
\end{tcolorbox}

\subsection*{ACTIVIDADES}

\begin{enumerate}[label=\arabic*., leftmargin=*]
    \item En la clase del 7 de Septiembre resolvimos analíticamente la EDO:
    \[ x'(t) = \mu - x^{2} \]
    para $\mu < 0$. Este es un modelo clásico de EDO autónoma con Bifurcación tipo Nodo-Silla, pues según el valor de $\mu$ las soluciones y el sistema presentan características diferentes.

    Ahora considera el caso $\mu > 0$. Responde:
    \begin{enumerate}[label=(\alph*), leftmargin=*]
        \item Da las soluciones de equilibrio.
        \subsection*{Solución}
        % Escribe tu respuesta aquí

        \item Resuelve analíticamente, dando los dominios de cada solución y región (imagen) comprendida.
        \subsection*{Solución}
        % Escribe tu respuesta aquí

        \item ¿Las soluciones que comienzan fuera de las de equilibrio, llegan a intersecar a alguna de las soluciones de equilibrio?
        \subsection*{Solución}
        % Escribe tu respuesta aquí

        \item Utilizando el software \texttt{Sol\_Analitica\_Bifurcacion\_Node-Saddle\_dx\_mu-x2\_mu\_neg.py}, modifícalo para obtener:
        \begin{itemize}[leftmargin=*]
            \item La solución sin condiciones iniciales (ic).
            \item Una solución con (ic), para cada una de las regiones en que separaron las soluciones de equilibrio el espacio.
        \end{itemize}
        \textbf{NOTA:} Eres libre de escoger las condiciones iniciales.
        \subsection*{Solución}
        % Escribe tu respuesta aquí

        \item La solución que obtuviste con el programa en el caso ``sin condiciones iniciales'', ¿a cuál equivale de las obtenidas analíticamente?
        \subsection*{Solución}
        % Escribe tu respuesta aquí
    \end{enumerate}

    \item \textbf{Control de congestión TCP (Modelo de fluidos)}

    En una red TCP, el tamaño de la ventana de congestión $W(t)$ crece linealmente y se reduce a la mitad cuando hay pérdidas. El modelo de fluidos de Misra et al. aproxima este comportamiento mediante una EDO que describe la evolución promedio de $W$.

    \begin{itemize}[leftmargin=*]
        \item \textbf{Ecuación Modelo:}
        \[ \frac{dW}{dt} = \frac{1}{R} - \frac{W}{2R}p \]
        donde $W(t)$: tamaño de la ventana TCP; $R$: RTT\footnote{RTT en Redes e Internet es el \textit{Round-Trip Time}. Es una métrica para medir latencia y rendimiento de una conexión a internet o una red local. Mide el tiempo exacto que pasa desde que un usuario envía un paquete de datos (como hacer clic en un enlace o cargar una página) hasta que el servidor responde y la información regresa al equipo de origen.}; $p$: probabilidad de pérdida. Se asume $R$ y $p$ constantes.
        \item \textbf{Condición inicial:} $W(0) = W_{0}$.
    \end{itemize}

    Preguntas:
    \begin{enumerate}[label=(\alph*), leftmargin=*]
        \item Halla $W(t)$.
        \subsection*{Solución}
        % Escribe tu respuesta aquí

        \item Determine el valor de equilibrio $W^{*}$.
        \subsection*{Solución}
        % Escribe tu respuesta aquí
    \end{enumerate}

    \item \textbf{Dinámica de un buffer de almacenamiento}

    Un buffer recibe datos a una tasa $\lambda$ y los procesa a una tasa proporcional a su ocupación $\mu B$. Este modelo de flujo describe cómo varía la ocupación del buffer con el tiempo.

    \begin{itemize}[leftmargin=*]
        \item \textbf{Ecuación Modelo:}
        \[ \frac{dB}{dt} = \lambda - \mu B \]
        \item \textbf{Condición inicial:} $B(0) = B_{0}$.
    \end{itemize}

    Tomando los valores:
    \[ \lambda = 250\text{ GB/s}, \quad \mu = 0.5\text{ s}^{-1}, \quad B_{0} = 50\text{ GB} \]
    cuya interpretación es:

    \vspace{0.3cm}
    \begin{center}
    \begin{tabular}{c l c}
    \toprule
    \textbf{Símbolo} & \textbf{Significado} & \textbf{Unidades} \\
    \midrule
    $B(t)$ & Ocupación del buffer en el tiempo $t$ & Megabytes (MB) \\
    $\lambda$ & Tasa de llegada de datos & $\text{MB/s}$ \\
    $\mu$ & Tasa de servicio por unidad de ocupación & $1/\text{s}$ \\
    $B_{0}$ & Ocupación inicial del buffer & MB \\
    \bottomrule
    \end{tabular}
    \end{center}
    \vspace{0.3cm}

    \begin{itemize}[leftmargin=*]
        \item Llegan $250\text{ GB}$ por segundo al buffer.
        \item El buffer procesa datos a una tasa de $0.5$ veces su ocupación actual (si tiene $20\text{ GB}$, procesa $10\text{ GB/s}$).
        \item Inicialmente el buffer contiene $50\text{ GB}$.
    \end{itemize}

    Preguntas:
    \begin{enumerate}[label=(\alph*), leftmargin=*]
        \item Halle la solución general $B(t)$.
        \subsection*{Solución}
        % Escribe tu respuesta aquí

        \item Aplique la condición inicial $B(0) = 50\text{ GB}$ y obtenga la solución particular.
        \subsection*{Solución}
        % Escribe tu respuesta aquí

        \item Determine el valor de equilibrio $B^{*}$. ¿Qué representa físicamente?
        \subsection*{Solución}
        % Escribe tu respuesta aquí

        \item Calcule el tiempo que tarda el buffer en alcanzar el 90\% de su valor de equilibrio.
        \subsection*{Solución}
        % Escribe tu respuesta aquí

        \item \textbf{Condición de estabilidad:} ¿Para qué valores de $\lambda$ y $\mu$ el buffer se estabiliza? ¿Qué pasa si $\lambda$ es demasiado grande?
        \subsection*{Solución}
        % Escribe tu respuesta aquí

        \item \textbf{Interpretación física:} Si el buffer tiene un tamaño máximo $Q = 500\text{ GB}$, ¿se desborda con los valores dados? ¿En qué instante?
        \subsection*{Solución}
        % Escribe tu respuesta aquí
    \end{enumerate}

    \item \textbf{Crecimiento de una base de datos distribuida}

    Una base de datos distribuida crece conforme pasa el tiempo a la vez que elimina registros obsoletos en ella a lo largo de este. Una EDO muy simplificada que predice el cambio de tamaño en la base de datos a lo largo del tiempo se conoce como \textit{modelo logístico con cosecha} y se expresa mediante la siguiente EDO:

    \begin{itemize}[leftmargin=*]
        \item \textbf{Ecuación Modelo:}
        \begin{equation*}
            \frac{dD}{dt} = rD\left(1 - \frac{D}{K}\right) - \delta D \tag{2}
        \end{equation*}
    \end{itemize}

    donde:

    \vspace{0.3cm}
    \begin{center}
    \resizebox{\linewidth}{!}{
    \begin{tabular}{c p{6.5cm} p{6.5cm}}
    \toprule
    \textbf{Variable} & \textbf{Significado en el modelo original (ecología)} & \textbf{Interpretación en el contexto ``base de datos''} \\
    \midrule
    $D(t)$ & Población (biomasa, número de individuos) en el tiempo $t$ & Tamaño de la base de datos (número de registros o GB) en el tiempo $t$ \\
    \addlinespace
    $r$ & Tasa de crecimiento intrínseco (natalidad menos mortalidad natural) & Tasa de crecimiento de la base de datos (inserciones por unidad de tiempo por registro existente) \\
    \addlinespace
    $K$ & Capacidad de carga del entorno (máximo sostenible) & Capacidad máxima de la base de datos (límite de almacenamiento o de rendimiento) \\
    \addlinespace
    $\delta$ & Tasa de cosecha (proporcional a la población) & Tasa de eliminación de registros obsoletos (proporcional al tamaño actual) \\
    \addlinespace
    $D_{0}$ & Población inicial al tiempo $t_{0}$ (usualmente $t_{0}=0$) & Tamaño inicial de la base de datos al tiempo $t_{0}$ \\
    \bottomrule
    \end{tabular}
    }
    \end{center}
    \vspace{0.3cm}

    \textbf{Valores con sentido aplicable en la realidad:} \\
    Para una base de datos distribuida de tamaño medio (por ejemplo, un sistema de logs o un almacén de datos operacional):

    \vspace{0.3cm}
    \begin{center}
    \resizebox{\linewidth}{!}{
    \begin{tabular}{c c c p{7cm}}
    \toprule
    \textbf{Parámetro} & \textbf{Valor} & \textbf{Unidades} & \textbf{Interpretación} \\
    \midrule
    $D_{0}$ ($t_{0}=0$) & $100$ & TB & Tamaño inicial de la base de datos \\
    \addlinespace
    $r$ & $0.05$ & $\text{TB}/(\text{TB}\cdot\text{día}) = 1/\text{día}$ & Tasa de crecimiento: 5\% del tamaño actual por día \\
    \addlinespace
    $K$ & $1000$ & TB & Capacidad máxima del clúster ($1\text{ PB}$) \\
    \addlinespace
    $\delta$ & $0.02$ & $1/\text{día}$ & Tasa de eliminación: 2\% del tamaño actual por día \\
    \bottomrule
    \end{tabular}
    }
    \end{center}
    \vspace{0.3cm}

    Preguntas:
    \begin{enumerate}[label=(\alph*), leftmargin=*]
        \item Muestra que la ED (2) se puede reescribir como:
        \[ \frac{dD}{dt} = (r-\delta)D\left(1 - \frac{D}{K(1-\delta/r)}\right) \]
        donde queda claro que la eliminación proporcional $\delta$ reduce tanto la tasa de crecimiento $r$ como la capacidad de carga $K$ efectivas.
        \subsection*{Solución}
        % Escribe tu respuesta aquí

        \item Calcula la tasa de crecimiento efectiva: $r-\delta$.
        \subsection*{Solución}
        % Escribe tu respuesta aquí

        \item Da la capacidad de carga efectiva: $K_{\text{eff}} = K(1-\delta/r)$.
        \subsection*{Solución}
        % Escribe tu respuesta aquí

        \item Determine el tamaño (solución) de equilibrio.
        \subsection*{Solución}
        % Escribe tu respuesta aquí

        \item ¿En cuántas regiones queda dividido el espacio, según las soluciones de equilibrio?
        \subsection*{Solución}
        % Escribe tu respuesta aquí

        \item ¿En qué región tiene sentido físico (real) nuestro problema?
        \subsection*{Solución}
        % Escribe tu respuesta aquí

        \item Halle $D(t)$, con la condición inicial y valores de los parámetros dados. \\
        \textbf{HINT:} Transforma la ED (2) a la forma:
        \[ \frac{dD}{dt} = \left[(r-\delta) - \frac{r}{K}D\right]D \]
        separa y resuelve el lado en términos de ``$D$'' por fracciones parciales.
        \subsection*{Solución}
        % Escribe tu respuesta aquí
    \end{enumerate}
\end{enumerate}